% !PRINTED 167
\weberpage{167}
Aus (2), (3), (4) ergiebt sich aber, wenn wir noch
\[ \Psi = \psi - a_1\psi_1 - a_2\psi_2 \cdots - a_\mu\psi_\mu \]
setzen:
\weq{5}{F = a_1F_1+a_2F_2+\cdots+a_\mu F_\mu+\Psi,}
worin $\Psi$ die Variable $y$ h\"ochstens bis zum Grade $r-1$ enth\"alt. Alle Functionen $\Psi$, die der Bedingung (5) gen\"ugen, bilden ein System $\mathfrak S_{r-1}$.

Wir machen nun weiter die Annahme, das Theorem I. sei bewiesen f\"ur Functionen eines jeden Systemes $\mathfrak S_{r-1}$. Dann k\"onnen wir also $\Psi$ durch eine Basis darstellen in der Form:
\weq{6}{\Psi = b_1\Psi_1+b_2\Psi_2+\cdots+b_\nu\Psi_\nu,}
worin $\Psi_1,\Psi_2,\ldots,\Psi_\nu$ specielle Functionen unseres Systemes $\mathfrak S_{r-1}$ sind. Man kann also \"uber die Co\"efficienten $a_{i,k}$ so verf\"ugen, dass die Functionen
\weq{7}{\begin{aligned}
F_1' &= a_{1,1}F_1+a_{2,1}F_2+\cdots+a_{\mu,1}F_\mu+\Psi_1\\
F_2' &= a_{1,2}F_1+a_{2,2}F_2+\cdots+a_{\mu,2}F_\mu+\Psi_2\\
&\cdots\cdots\cdots\cdots\cdots\cdots\cdots\cdots\cdots\cdots\cdots\cdots\cdots\cdots\\
F_\nu' &= a_{1,\nu}F_1+a_{2,\nu}F_2+\cdots+a_{\mu,\nu}F_\mu+\Psi_\nu
\end{aligned}}
in $\mathfrak S_r$ enthalten sind.

Multiplicirt man mit $b_1,b_2,\ldots b_\nu$ und setzt in der Summe nach (5) und (6):
\[ \begin{aligned}
b_1\Psi_1+b_2\Psi_2+\cdots+b_\nu\Psi_\nu &= \Psi\\
&= F - a_1F_1-a_2F_2-\cdots-a_\mu F_\mu,
\end{aligned} \]
so folgt
\weq{8}{\begin{aligned}
F &= b_1F_1'+b_2F_2'+\cdots+b_\nu F_\nu'\\
&\quad+A_1F_1+A_2F_2+\cdots+A_\mu F_\mu,
\end{aligned}}
worin
\weq{9}{\begin{aligned}
A_1 &= a_1 - a_{1,1}b_1 - \cdots - a_{1,\nu}b_\nu\\
&\cdots\cdots\cdots\cdots\cdots\cdots\cdots\cdots\cdots\cdots\cdots\\
A_\mu &= a_\mu - a_{\mu,1}b_1 - \cdots - a_{\mu,\nu}b_\nu.
\end{aligned}}
Die Formen $F_1',\ldots F_\nu', F_1,\ldots F_\mu$ geh\"oren dem Systeme $\mathfrak S_r$ an und bilden eine Basis dieses Systemes, und das Theorem I. ist also unter der Voraussetzung, dass es f\"ur $\mathfrak S_0$ gilt, f\"ur jedes System $\mathfrak S_r$ bewiesen.

Es sei nun $\mathfrak S$ ein beliebiges System von Formen von $n+1$ Variablen. Wir greifen irgend eine Form $F_0$ vom Grade $r$ aus diesem Systeme heraus, und setzen f\"ur die Variablen, von denen das System abh\"angt,
\weq{10}{y,\ x_1+\lambda_1y,\ \ldots,\ x_n+\lambda_ny,}
%CONT>>
